\chapter{2005年江苏卷}
\section{单选题}
\begin{problem}
设集合 \(A = \{1, 2\}\)，\(B = \{1, 2, 3\}\)，\(C = \{2, 3, 4\}\)，则 \(\left(A \cap B\right) \cup C =\)（\quad）

\choices
    {\(\{1,2,3\}\)}
    {\(\{1,2,4\}\)}
    {\(\{2,3,4\}\)}
    {\(\{1,2,3,4\}\)}
\end{problem}

\begin{answer}
D
\end{answer}

\begin{solution}
由 \(A=\{1,2\}\)，\(B=\{1,2,3\}\) 可知 \(A\subset B\)，因此 \(A\cap B=A=\{1,2\}\). 再与 \(C=\{2,3,4\}\) 求并集，得到

\(\left(A\cap B\right)\cup C=\{1,2\}\cup\{2,3,4\}=\{1,2,3,4\}\).

所以应选 D.
\end{solution}

\begin{problem}
函数 \(y = 2^{1 - x} + 3\)（\(x \in \R\)）的反函数的解析表达式为（\quad）

\choices
    {\(y=\log_{2}\frac{2}{x - 3}\)}
    {\(y=\log_{2}\frac{x - 3}{2}\)}
    {\(y=\log_{2}\frac{3 - x}{2}\)}
    {\(y=\log_{2}\frac{2}{3 - x}\)}
\end{problem}

\begin{answer}
A
\end{answer}

\begin{solution}
由 \(y=2^{1-x}+3\)，先移项得 \(y-3=2^{1-x}>0\)，所以原函数的值域为 \(y>3\). 对等式两边取以 \(2\) 为底的对数，得

\(\log_{2}(y-3)=1-x\)，从而 \(x=1-\log_{2}(y-3)=\log_{2}\dfrac{2}{y-3}\).

交换自变量与因变量的记号，并由反函数定义域为 \(x>3\)，得到反函数 \(y=\log_{2}\dfrac{2}{x-3}\). 所以应选 A.
\end{solution}

\begin{problem}
在各项都为正数的等比数列 \(\{a_{n}\}\) 中，首项 \(a_1 = 3\)，前三项和为 \(21\)，则 \(a_3 + a_4 + a_5 =\)（\quad）

\choices
    {\(33\)}
    {\(72\)}
    {\(84\)}
    {\(189\)}
\end{problem}

\begin{answer}
C
\end{answer}

\begin{solution}
设等比数列的公比为 \(q\). 由 \(a_1=3\) 和前三项的和为 \(21\)，有

\(3+3q+3q^2=21\)，即 \(q^2+q-6=0\)，所以 \(q=2\) 或 \(q=-3\).

因为数列各项均为正数，公比不能为负数，故 \(q=2\). 于是 \(a_3=3q^2=12\)，\(a_4=3q^3=24\)，\(a_5=3q^4=48\)，从而

\(a_3+a_4+a_5=12+24+48=84\).

所以应选 C.
\end{solution}

\begin{problem}
在正三棱柱 \(ABC-A_{1}B_{1}C_{1}\) 中，若 \(AB = 2\)，\(AA_{1} = 1\)，则点 \(A\) 到平面 \(A_{1}BC\) 的距离为（\quad）

\choices
    {\(\frac{\sqrt{3}}{4}\)}
    {\(\frac{\sqrt{3}}{2}\)}
    {\(\frac{3\sqrt{3}}{4}\)}
    {\(\sqrt{3}\)}
\end{problem}

\begin{answer}
B
\end{answer}

\begin{solution}
建立直角坐标系，使 \(A=(0,0,0)\)，\(B=(2,0,0)\)，\(C=(1,\sqrt{3},0)\)，且棱柱的高沿 \(z\) 轴，故 \(A_1=(0,0,1)\). 于是

\(\overrightarrow{A_1B}=(2,0,-1)\)，\(\overrightarrow{A_1C}=(1,\sqrt{3},-1)\).

平面 \(A_1BC\) 的一个法向量为

\(\overrightarrow{A_1B}\times\overrightarrow{A_1C}=(\sqrt{3},1,2\sqrt{3})\).

又 \(\overrightarrow{A_1A}=(0,0,-1)\)，所以点 \(A\) 到平面 \(A_1BC\) 的距离为

\(\dfrac{\left|\overrightarrow{A_1A}\cdot(\sqrt{3},1,2\sqrt{3})\right|}{\sqrt{(\sqrt{3})^2+1^2+(2\sqrt{3})^2}}=\dfrac{2\sqrt{3}}{4}=\dfrac{\sqrt{3}}{2}\).

所以应选 B.
\end{solution}

\begin{problem}
\(\triangle ABC\) 中，\(\angle A = \frac{\pi}{3}\)，\(BC = 3\)，则 \(\triangle ABC\) 的周长为（\quad）

\choices
    {\(4\sqrt{3}\sin \left(B + \frac{\pi}{3}\right) + 3\)}
    {\(4\sqrt{3}\sin \left(B + \frac{\pi}{6}\right) + 3\)}
    {\(6\sin \left(B + \frac{\pi}{3}\right) + 3\)}
    {\(6\sin \left(B + \frac{\pi}{6}\right) + 3\)}
\end{problem}

\begin{answer}
D
\end{answer}

\begin{solution}
在 \(\triangle ABC\) 中，令 \(a=BC=3\)，\(b=CA\)，\(c=AB\). 由正弦定理，

\(\dfrac{b}{\sin B}=\dfrac{a}{\sin A}=\dfrac{3}{\sin(\pi/3)}\)，所以 \(b=2\sqrt{3}\sin B\).

又 \(A+B+C=\pi\)，故 \(C=\dfrac{2\pi}{3}-B\). 同理，

\(c=\dfrac{3\sin C}{\sin(\pi/3)}=2\sqrt{3}\sin\left(\dfrac{2\pi}{3}-B\right)=2\sqrt{3}\sin\left(B+\dfrac{\pi}{3}\right)\).

因此周长为

\(a+b+c=3+2\sqrt{3}\left[\sin B+\sin\left(B+\dfrac{\pi}{3}\right)\right]\).

利用和差化积公式，

\(\sin B+\sin\left(B+\dfrac{\pi}{3}\right)=2\sin\left(B+\dfrac{\pi}{6}\right)\cos\dfrac{\pi}{6}=\sqrt{3}\sin\left(B+\dfrac{\pi}{6}\right)\).

所以周长为 \(6\sin\left(B+\dfrac{\pi}{6}\right)+3\)，应选 D.
\end{solution}

\begin{problem}
抛物线 \(y = 4x^{2}\) 上的一点 \(M\) 到焦点的距离为 \(1\)，则点 \(M\) 的纵坐标是（\quad）

\choices
    {\(\frac{17}{16}\)}
    {\(\frac{15}{16}\)}
    {\(\frac{7}{8}\)}
    {\(0\)}
\end{problem}

\begin{answer}
B
\end{answer}

\begin{solution}
将抛物线写成标准形式 \(x^2=4py\). 由 \(y=4x^2\) 得 \(x^2=\dfrac14y\)，所以 \(4p=\dfrac14\)，即 \(p=\dfrac1{16}\). 因此焦点为 \(F\left(0,\dfrac1{16}\right)\)，准线为 \(y=-\dfrac1{16}\).

设 \(M=(x,y)\). 由抛物线的定义，\(MF\) 等于点 \(M\) 到准线的距离. 由于 \(y=4x^2\geqslant0\)，点 \(M\) 到准线的距离为 \(y+\dfrac1{16}\). 题设 \(MF=1\)，故

\(1=y+\dfrac1{16}\)，从而 \(y=\dfrac{15}{16}\).

所以应选 B.
\end{solution}

\begin{problem}
在一次歌手大奖赛上，七位评委为歌手打出的分数如下：\(9.4\)，\(8.4\)，\(9.4\)，\(9.9\)，\(9.6\)，\(9.4\)，\(9.7\). 去掉一个最高分和一个最低分后，所剩数据的平均值和方差分别为（\quad）

\choices
    {\(9.4, 0.484\)}
    {\(9.4, 0.016\)}
    {\(9.5, 0.04\)}
    {\(9.5, 0.016\)}
\end{problem}

\begin{answer}
D
\end{answer}

\begin{solution}
最高分为 \(9.9\)，最低分为 \(8.4\). 去掉这两项后，剩余数据为 \(9.4\)，\(9.4\)，\(9.4\)，\(9.6\)，\(9.7\)，共有 \(5\) 个数.

它们的平均值为

\(\overline{x}=\dfrac{3\times9.4+9.6+9.7}{5}=\dfrac{47.5}{5}=9.5\).

按总体方差的定义，方差为

\(\dfrac{3(9.4-9.5)^2+(9.6-9.5)^2+(9.7-9.5)^2}{5}=\dfrac{3\times0.01+0.01+0.04}{5}=0.016\).

所以应选 D.
\end{solution}

\begin{problem}
设 \(\alpha\)，\(\beta\)，\(\gamma\) 为两两不重合的平面，\(l\)，\(m\)，\(n\) 为两两不重合的直线，给出下列四个命题：
\begin{circlelist}
\item 若 \(\alpha \perp \gamma\)，\(\beta \perp \gamma\)，则 \(\alpha \parallel \beta\)；
\item 若 \(m \subset \alpha\)，\(n \subset \alpha\)，\(m \parallel \beta\)，\(n \parallel \beta\)，则 \(\alpha \parallel \beta\)；
\item 若 \(\alpha \parallel \beta\)，\(l \subset \alpha\)，则 \(l \parallel \beta\)；
\item 若 \(\alpha \cap \beta = l\)，\(\beta \cap \gamma = m\)，\(\gamma \cap \alpha = n\)，\(l \parallel \gamma\)，则 \(m \parallel n\).
\end{circlelist}
其中真命题的个数是（\quad）

\choices
    {\(1\)}
    {\(2\)}
    {\(3\)}
    {\(4\)}
\end{problem}

\begin{answer}
B
\end{answer}

\begin{solution}
命题 \(\circled{1}\) 不成立. 例如令 \(\gamma\) 为水平面，令 \(\alpha\)，\(\beta\) 为两个相交的竖直平面. 两个竖直平面都垂直于水平面，但它们彼此相交，不能推出 \(\alpha\parallel\beta\).

命题 \(\circled{2}\) 不成立. 取 \(\alpha\) 为平面 \(z=0\)，\(\beta\) 为平面 \(y=0\). 在 \(\alpha\) 内取直线 \(m\colon y=1\)，\(z=0\) 和 \(n\colon y=2\)，\(z=0\)，则 \(m\)，\(n\) 都与 \(\beta\) 平行，而 \(\alpha\) 与 \(\beta\) 相交于直线 \(y=0\)，\(z=0\)，所以不能推出 \(\alpha\parallel\beta\).

命题 \(\circled{3}\) 成立. 若 \(\alpha\parallel\beta\)，则 \(\alpha\) 内任一直线都不可能与 \(\beta\) 相交. 因为 \(l\subset\alpha\)，所以 \(l\parallel\beta\).

命题 \(\circled{4}\) 成立. 由 \(l\subset\alpha\)，\(l\parallel\gamma\)，且 \(n=\alpha\cap\gamma\)，可知 \(l\parallel n\). 同理，由 \(l\subset\beta\)，\(l\parallel\gamma\)，且 \(m=\beta\cap\gamma\)，可知 \(l\parallel m\). 两条都与 \(l\) 平行的直线互相平行，因此 \(m\parallel n\).

只有命题 \(\circled{3}\)，\(\circled{4}\) 正确，真命题有 \(2\) 个，所以应选 B.
\end{solution}

\begin{problem}
设 \(k = 1\)，\(2\)，\(3\)，\(4\)，\(5\)，则 \((x + 2)^5\) 的展开式中 \(x^k\) 的系数不可能是（\quad）

\choices
    {\(10\)}
    {\(40\)}
    {\(50\)}
    {\(80\)}
\end{problem}

\begin{answer}
C
\end{answer}

\begin{solution}
由二项式定理，

\((x+2)^5=\displaystyle\sum_{k=0}^{5}\mathrm C_{5}^{k}x^k2^{5-k}\).

因此 \(x^k\) 的系数为 \(\mathrm C_{5}^{k}2^{5-k}\). 分别代入 \(k=1\)，\(2\)，\(3\)，\(4\)，\(5\)，所得系数为

\(\mathrm C_{5}^{1}2^4=80\)，\(\mathrm C_{5}^{2}2^3=80\)，\(\mathrm C_{5}^{3}2^2=40\)，\(\mathrm C_{5}^{4}2=10\)，\(\mathrm C_{5}^{5}=1\).

选项中的 \(10\)，\(40\)，\(80\) 均出现，而 \(50\) 没有出现，所以应选 C.
\end{solution}

\begin{problem}
若 \(\sin \left(\frac{\pi}{6} - \alpha\right) = \frac{1}{3}\)，则 \(\cos \left(\frac{2\pi}{3} + 2\alpha\right) =\)（\quad）

\choices
    {\(\displaystyle -\frac{7}{9}\)}
    {\(\displaystyle -\frac{1}{3}\)}
    {\(\displaystyle \frac{1}{3}\)}
    {\(\displaystyle \frac{7}{9}\)}
\end{problem}

\begin{answer}
A
\end{answer}

\begin{solution}
令 \(t=\dfrac{\pi}{6}-\alpha\)，则 \(\sin t=\dfrac13\)，且 \(\alpha=\dfrac{\pi}{6}-t\). 因而

\(\dfrac{2\pi}{3}+2\alpha=\dfrac{2\pi}{3}+2\left(\dfrac{\pi}{6}-t\right)=\pi-2t\).

利用 \(\cos(\pi-u)=-\cos u\) 和二倍角公式，得到

\(\cos\left(\dfrac{2\pi}{3}+2\alpha\right)=\cos(\pi-2t)=-\cos2t=-\left(1-2\sin^2t\right)=-1+\dfrac{2}{9}=-\dfrac79\).

所以应选 A.
\end{solution}

\begin{problem}
点 \(P(-3,1)\) 在椭圆 \(\frac{x^2}{a^2} + \frac{y^2}{b^2} = 1\)（\(a > b > 0\)）的左准线上. 过点 \(P\) 且方向为 \(\bs{a} = (2,-5)\) 的光线，经直线 \(y = -2\) 反射后通过椭圆的左焦点，则这个椭圆的离心率为（\quad）

\choices
    {\(\frac{\sqrt{3}}{3}\)}
    {\(\frac{1}{3}\)}
    {\(\frac{\sqrt{2}}{2}\)}
    {\(\frac{1}{2}\)}
\end{problem}

\begin{answer}
A
\end{answer}

\begin{solution}
设椭圆的半长轴为 \(a\)，焦距的一半为 \(c\)，离心率为 \(e=\dfrac ca\). 左准线方程为 \(x=-\dfrac{a}{e}\). 由 \(P(-3,1)\) 在左准线上，得

\(\dfrac{a}{e}=3\)，即 \(a=3e\)，从而 \(c=ae=3e^2\).

入射光线从 \(P\) 沿方向向量 \((2,-5)\) 前进，可写为 \(P+t(2,-5)\). 与 \(y=-2\) 相交时，

\(1-5t=-2\)，所以 \(t=\dfrac35\)，交点为

\(Q=(-3,1)+\dfrac35(2,-5)=\left(-\dfrac95,-2\right)\).

水平直线反射保持横向分量不变、使竖直分量变号，因此反射光线的方向向量为 \((2,5)\). 从 \(Q\) 到 \(y=0\) 需要参数 \(s=\dfrac25\)，对应点的横坐标为

\(-\dfrac95+2\times\dfrac25=-1\).

所以反射光线经过 \((-1,0)\). 该点是椭圆左焦点 \((-c,0)\)，故 \(c=1\). 代入 \(c=3e^2\)，得 \(3e^2=1\)，又 \(e>0\)，所以 \(e=\dfrac{\sqrt{3}}{3}\).

因此应选 A.
\end{solution}

\begin{problem}
四棱锥的 \(8\) 条棱代表 \(8\) 种不同的化工产品，有公共点的两条棱代表的化工产品放在同一仓库是危险的，没有公共顶点的两条棱所代表的化工产品放在同一仓库是安全的，现打算用编号为 \(\circled{1}\)，\(\circled{2}\)，\(\circled{3}\)，\(\circled{4}\) 的 \(4\) 个仓库存放这 \(8\) 种化工产品，那么安全存放的不同方法种数为（\quad）

\choices
    {\(96\)}
    {\(48\)}
    {\(24\)}
    {\(0\)}
\end{problem}

\begin{answer}
B
\end{answer}

\begin{solution}
设四棱锥底面为 \(ABCD\)，顶点为 \(S\). 同一仓库中的任意两条棱都没有公共顶点；若同一仓库中有三条棱，则这三条棱需要六个互不相同的顶点，但四棱锥只有五个顶点，所以每个仓库至多放两条棱. 四条侧棱 \(SA\)，\(SB\)，\(SC\)，\(SD\) 两两有公共顶点，因此同一仓库中至多放一条侧棱. 另一方面，底面四条棱构成一个四边形，能同仓存放的两条底面棱只能是对边.

共有 \(8\) 条棱和 \(4\) 个仓库. 由于每个仓库最多容纳两条棱，安全存放时每个仓库恰好容纳两条棱；又四条侧棱必须分别进入四个不同仓库，所以每个仓库恰好含一条侧棱和一条底面棱.

两条棱没有公共顶点时才可同仓，因而底面棱的安全配对为

\(AB\) 可与 \(SC\)，\(SD\) 配对，\(BC\) 可与 \(SD\)，\(SA\) 配对，\(CD\) 可与 \(SA\)，\(SB\) 配对，\(DA\) 可与 \(SB\)，\(SC\) 配对.

要使四条底面棱分别与四条侧棱配成安全对，只有两种完全配法：

\(AB\text{-}SC\)，\(\ BC\text{-}SD\)，\(\ CD\text{-}SA\)，\(\ DA\text{-}SB\)，或 \(AB\text{-}SD\)，\(\ BC\text{-}SA\)，\(\ CD\text{-}SB\)，\(\ DA\text{-}SC\).

每一种配法得到的四个安全对还可以任意放入四个有编号的仓库，有 \(4!\) 种. 所以总方法数为 \(2\times4!=48\)，应选 B.
\end{solution}

\section{填空题}
\begin{problem}
命题“若 \(a > b\)，则 \(2^a > 2^b - 1\)”的否命题为\fillinblank{}.
\end{problem}

\begin{answer}
若 \(a\leqslant b\)，则 \(2^a\leqslant2^b-1\)
\end{answer}

\begin{solution}
令 \(p\) 表示条件 \(a>b\)，令 \(q\) 表示结论 \(2^a>2^b-1\). 原命题是“若 \(p\)，则 \(q\)”的形式.

否命题的形式是“若 \(\lnot p\)，则 \(\lnot q\)”. 因为 \(\lnot p\) 为 \(a\leqslant b\)，而严格不等式 \(2^a>2^b-1\) 的否定为 \(2^a\leqslant2^b-1\)，所以否命题为“若 \(a\leqslant b\)，则 \(2^a\leqslant2^b-1\)”.
\end{solution}

\begin{problem}
曲线 \(y = x^{3} + x + 1\) 在点 \((1,3)\) 处的切线方程是\fillinblank{}.
\end{problem}

\begin{answer}
\(4x-y-1=0\)
\end{answer}

\begin{solution}
设 \(f(x)=x^3+x+1\). 则 \(f'(x)=3x^2+1\)，所以在 \(x=1\) 处的切线斜率为 \(f'(1)=4\). 题设点为 \((1,3)\)，由点斜式得切线方程

\(y-3=4(x-1)\).

整理为 \(4x-y-1=0\).
\end{solution}

\begin{problem}
函数 \(y = \sqrt{\log_{0.5}(4x^2 - 3x)}\) 的定义域为\fillinblank{}.
\end{problem}

\begin{answer}
\(\left[-\frac{1}{4},0\right)\cup\left(\frac{3}{4},1\right]\)
\end{answer}

\begin{solution}
要使根式有意义，必须满足

\(\log_{0.5}(4x^2-3x)\geqslant0\).

由于对数真数必须为正，且底数 \(0.5\in(0,1)\)，函数 \(\log_{0.5}u\) 在 \(u>0\) 时递减，并且 \(\log_{0.5}u\geqslant0\) 等价于 \(0<u\leqslant1\). 所以需要同时满足

\(0<4x^2-3x\leqslant1\).

其中，\(4x^2-3x=x(4x-3)>0\) 给出 \(x<0\) 或 \(x>\dfrac34\). 另一方面，

\(4x^2-3x\leqslant1\iff4x^2-3x-1\leqslant0\iff(4x+1)(x-1)\leqslant0\)，

所以 \(-\dfrac14\leqslant x\leqslant1\). 取两组条件的交集，定义域为 \(\left[-\dfrac14,0\right)\cup\left(\dfrac34,1\right]\).
\end{solution}

\begin{problem}
若 \(3^{a} = 0.618\)，\(a \in [k,k+1), k \in \Z\)，则 \(k =\)\fillinblank{}.
\end{problem}

\begin{answer}
\(-1\)
\end{answer}

\begin{solution}
因为 \(3^x\) 的底数大于 \(1\)，所以它在 \(\R\) 上严格递增. 又

\(3^{-1}=\dfrac13<0.618<1=3^0\)，

从而 \(-1<a<0\). 这个开区间包含在且只包含在半开区间 \([-1,0)\) 中. 因此题设中的整数 \(k\) 为 \(-1\).
\end{solution}

\begin{problem}
已知 \(a\)，\(b\) 为常数，若 \(f(x) = x^{2} + 4x + 3\)，\(f(ax + b) = x^{2} + 10x + 24\)，则 \(5a - b =\)\fillinblank{}.
\end{problem}

\begin{answer}
\(2\)
\end{answer}

\begin{solution}
将 \(ax+b\) 代入 \(f(x)=x^2+4x+3\)，得

\(f(ax+b)=(ax+b)^2+4(ax+b)+3=a^2x^2+(2ab+4a)x+(b^2+4b+3)\).

与 \(x^2+10x+24\) 比较同次幂系数，得到

\(a^2=1\)，\(2a(b+2)=10\)，\(b^2+4b+3=24\).

由 \(a^2=1\)，分两种情况：

当 \(a=1\) 时，\(2(b+2)=10\)，所以 \(b=3\)，且 \(b^2+4b+3=24\) 成立；

当 \(a=-1\) 时，\(-2(b+2)=10\)，所以 \(b=-7\)，且 \(b^2+4b+3=24\) 也成立.

两种情形中，\(5a-b\) 分别为 \(5-3=2\) 和 \(-5-(-7)=2\). 所以所求值为 \(2\).
\end{solution}

\begin{problem}
在 \(\triangle ABC\) 中，\(O\) 为中线 \(AM\) 上的一个动点，若 \(AM = 2\)，则 \(\overrightarrow{OA} \cdot (\overrightarrow{OB} + \overrightarrow{OC})\) 的最小值是\fillinblank{}.
\end{problem}

\begin{answer}
\(-2\)
\end{answer}

\begin{solution}
因为 \(M\) 是边 \(BC\) 的中点，取 \(M\) 为原点，设 \(\overrightarrow{MA}=\bs{u}\)，\(\overrightarrow{MB}=\bs{v}\)，则 \(\overrightarrow{MC}=-\bs{v}\)，且 \(|\bs{u}|=AM=2\).

设动点 \(O\) 在线段 \(AM\) 上，令 \(\overrightarrow{MO}=t\bs{u}\)，其中 \(0\leqslant t\leqslant1\). 则

\(\overrightarrow{OA}=(1-t)\bs{u}\)，\(\overrightarrow{OB}=\bs{v}-t\bs{u}\)，\(\overrightarrow{OC}=-\bs{v}-t\bs{u}\).

因此 \(\overrightarrow{OB}+\overrightarrow{OC}=-2t\bs{u}\)，所求数量为

\(\overrightarrow{OA}\cdot(\overrightarrow{OB}+\overrightarrow{OC})=(1-t)\bs{u}\cdot(-2t\bs{u})=-2t(1-t)|\bs{u}|^2=-8t(1-t)\).

因为 \(0\leqslant t\leqslant1\)，有 \(t(1-t)\leqslant\dfrac14\)，且等号在 \(t=\dfrac12\) 时成立. 所以所求最小值为 \(-8\times\dfrac14=-2\).
\end{solution}

\section{解答题}
\begin{problem}
如图，圆 \(O_{1}\) 与圆 \(O_{2}\) 的半径都是 \(1\)，\(O_{1}O_{2} = 4\)，过动点 \(P\) 分别作圆 \(O_{1}\)，圆 \(O_{2}\) 的切线 \(PM\)，\(PN\)（\(M\)，\(N\) 分别为切点），使得 \(PM = \sqrt{2}PN\). 试建立适当的坐标系，并求动点 \(P\) 的轨迹方程.

\FigureLayoutDeclare{img/2005/jiangsu/q19_fig1.png}{7.2cm}{12bp 38bp 72bp 15bp}
\bitmapfigure[max width=0.42\linewidth]{img/2005/jiangsu/q19_fig1.png}
\end{problem}

\begin{solution}
取 \(O_1O_2\) 的中点 \(O\) 为原点，令 \(O_1O_2\) 所在直线为 \(x\) 轴，建立平面直角坐标系. 因为 \(O_1O_2=4\)，所以 \(O_1=(-2,0)\)，\(O_2=(2,0)\). 设动点 \(P=(x,y)\).

由切线长定理，

\(PM^2=PO_1^2-1=(x+2)^2+y^2-1\)，

\(PN^2=PO_2^2-1=(x-2)^2+y^2-1\).

题设 \(PM=\sqrt{2}PN\)，两边平方后得到

\((x+2)^2+y^2-1=2\left[(x-2)^2+y^2-1\right]\).

展开并移项，得 \(x^2-12x+y^2+3=0\)，配方为

\((x-6)^2+y^2=33\).

记该圆的圆心为 \(O'=(6,0)\). 反过来，若 \(P\) 在该圆上，则 \(PO'=\sqrt{33}\)，且 \(O'O_2=4\)，\(O'O_1=8\). 由三角不等式，

\(PO_2\geqslant O'O_2-PO'=\sqrt{33}-4>1\)，\(PO_1\geqslant O'O_1-PO'=8-\sqrt{33}>1\)，

所以 \(P\) 在两个半径为 \(1\) 的圆外，所需切线确实存在. 因此动点 \(P\) 的轨迹方程为 \((x-6)^2+y^2=33\).
\end{solution}

\begin{problem}
甲，乙两人各射击一次，击中目标的概率分别是 \(\frac{2}{3}\) 和 \(\frac{3}{4}\). 假设两人射击是否击中目标，相互之间没有影响；每次射击是否击中目标，相互之间没有影响.
\begin{enumerate}
\item 求甲射击 \(4\) 次，至少 \(1\) 次未击中目标的概率；
\item 求两人各射击 \(4\) 次，甲恰好击中目标 \(2\) 次且乙恰好击中目标 \(3\) 次的概率；
\item 假设某人连续 \(2\) 次未击中目标，则停止射击. 问：乙恰好射击 \(5\) 次后，被中止射击的概率是多少？
\end{enumerate}
\end{problem}

\begin{solution}
\begin{enumerate}
\item 甲四次射击全部击中的概率为 \(\left(\dfrac23\right)^4\). “至少一次未击中”是“全部击中”的对立事件，所以所求概率为

\(1-\left(\dfrac23\right)^4=1-\dfrac{16}{81}=\dfrac{65}{81}\).
\item 甲恰好击中 \(2\) 次的概率为

\(\mathrm C_4^2\left(\dfrac23\right)^2\left(1-\dfrac23\right)^2=6\times\dfrac49\times\dfrac19=\dfrac{8}{27}\).

乙恰好击中 \(3\) 次的概率为

\(\mathrm C_4^3\left(\dfrac34\right)^3\left(1-\dfrac34\right)=4\times\dfrac{27}{64}\times\dfrac14=\dfrac{27}{64}\).

甲与乙的射击结果相互独立，因此所求概率为

\(\dfrac{8}{27}\times\dfrac{27}{64}=\dfrac18\).
\item 乙恰好在第 \(5\) 次射击后被中止，必须满足第 \(4\)，\(5\) 次均未击中；为了不在第 \(4\) 次就停止，第 \(3\) 次必须击中；为了不在第 \(2\) 次就停止，前两次不能同时未击中. 这些条件合起来的概率为

\(\left[1-\left(\dfrac14\right)^2\right]\times\dfrac34\times\dfrac14\times\dfrac14=\dfrac{15}{16}\times\dfrac34\times\dfrac1{16}=\dfrac{45}{1024}\).
\end{enumerate}
\end{solution}

\begin{problem}
如图，在五棱锥 \(S-ABCDE\) 中，\(SA \perp\) 底面 \(ABCDE\)，\(SA = AB = AE = 2\)，\(BC = DE = \sqrt{3}\)，\(\angle BAE = \angle BCD = \angle CDE = 120^{\circ}\).
\begin{enumerate}
\item 求异面直线 \(CD\) 与 \(SB\) 所成的角（用反三角函数值表示）；
\item 证明：\(BC \perp\) 平面 \(SAB\)；
\item 用反三角函数值表示二面角 \(B-SC-D\) 的大小（本小问不必写出解答过程）.
\end{enumerate}

\bitmapfigure[max width=0.42\linewidth]{img/2005/jiangsu/q21_fig1.png}
\end{problem}

\begin{solution}
\begin{enumerate}
\item 延长 \(BC\)，\(ED\)，使它们相交于 \(F\). 因为 \(CF\) 与 \(CB\) 方向相反，\(DF\) 与 \(DE\) 方向相反，所以

\(\angle DCF=180^{\circ}-\angle BCD=60^{\circ}\)，\(\angle CDF=180^{\circ}-\angle CDE=60^{\circ}\).

于是 \(\triangle CDF\) 的两个角为 \(60^{\circ}\)，它是正三角形，故 \(CF=DF\). 又因为 \(B\)，\(C\)，\(F\) 与 \(E\)，\(D\)，\(F\) 分别共线，且 \(BC=DE\)，所以 \(BF=EF\). 此外，\(\angle BFE=\angle CFD=60^{\circ}\)，故 \(\triangle BFE\) 也是正三角形，从而 \(BF=EF=BE\).

在等腰三角形 \(ABE\) 中，\(AB=AE=2\)，且 \(\angle BAE=120^{\circ}\). 由余弦定理，

\(BE^2=AB^2+AE^2-2AB\cdot AE\cos120^{\circ}=4+4+4=12\)，

所以 \(BE=2\sqrt{3}\). 因而 \(BF=EF=BE=2\sqrt{3}\)，且 \(\triangle BFE\) 为正三角形. 于是 \(\angle FBE=\angle FCD=60^{\circ}\)，又 \(BF\) 与 \(CF\) 共线，所以 \(BE\parallel CD\).

因此异面直线 \(CD\) 与 \(SB\) 所成的角等于直线 \(BE\) 与 \(SB\) 所成的锐角 \(\angle SBE\). 因为 \(SA\perp\) 底面，\(SA=AB=AE=2\)，所以

\(SB=\sqrt{SA^2+AB^2}=2\sqrt{2}\)，\(SE=\sqrt{SA^2+AE^2}=2\sqrt{2}\).

在 \(\triangle SBE\) 中，\(BE=2\sqrt{3}\)，由余弦定理，

\(\cos\angle SBE=\dfrac{SB^2+BE^2-SE^2}{2SB\cdot BE}=\dfrac{8+12-8}{2\cdot2\sqrt{2}\cdot2\sqrt{3}}=\dfrac{\sqrt{6}}{4}\).

所以所求角为 \(\arccos\dfrac{\sqrt{6}}{4}\).
\item 由 \(AB=AE\) 和 \(\angle BAE=120^{\circ}\)，得等腰三角形 \(ABE\) 的底角 \(\angle ABE=30^{\circ}\). 又 \(BF\) 与 \(BC\) 同向，且 \(\triangle BFE\) 为正三角形，所以 \(\angle EBC=\angle FBE=60^{\circ}\). 因此

\(\angle ABC=\angle ABE+\angle EBC=30^{\circ}+60^{\circ}=90^{\circ}\)，

即 \(BC\perp AB\). 又因为 \(SA\perp\) 底面 \(ABCDE\)，而 \(BC\) 在底面内，所以 \(BC\perp SA\). 直线 \(AB\)，\(SA\) 是平面 \(SAB\) 内相交的两条直线，故 \(BC\perp\) 平面 \(SAB\).
\item 令 \(r=\sqrt{3}\)，取底面为 \(z=0\) 平面，建立直角坐标系

\(A=(0,0,0)\)，\(S=(0,0,2)\)，\(B=(2,0,0)\)，\(C=(2,r,0)\)，\(D=\left(\dfrac12,\dfrac{3r}{2},0\right)\)，\(E=(-1,r,0)\).

这些坐标满足 \(AB=AE=2\)，\(BC=DE=r\)，以及三个给定的 \(120^{\circ}\) 角，因此可用于计算二面角. 在平面 \(BSC\) 和平面 \(DSC\) 中，分别取相交向量 \(\overrightarrow{SC}\)，\(\overrightarrow{SB}\) 与 \(\overrightarrow{SC}\)，\(\overrightarrow{SD}\)，则两平面的法向量可分别取为

\(\bs{n}_1=\overrightarrow{SC}\times\overrightarrow{SB}=(-2r,0,-2r)\)，

\(\bs{n}_2=\overrightarrow{SC}\times\overrightarrow{SD}=\left(r,3,\dfrac{5r}{2}\right)\).

这两个法向量分别对应二面角两侧指向 \(B\)，\(D\) 的半平面，所以它们的夹角就是所求二面角 \(\theta\). 计算得

\(\bs{n}_1\cdot\bs{n}_2=-21\)，\(|\bs{n}_1|=2\sqrt{6}\)，\(|\bs{n}_2|=\dfrac{\sqrt{123}}{2}\).

于是

\(\cos\theta=\dfrac{\bs{n}_1\cdot\bs{n}_2}{|\bs{n}_1||\bs{n}_2|}=-\dfrac{21}{3\sqrt{82}}=-\dfrac{7\sqrt{82}}{82}\).

所以 \(\theta=\arccos\left(-\dfrac{7\sqrt{82}}{82}\right)=\pi-\arccos\dfrac{7\sqrt{82}}{82}\).
\end{enumerate}
\end{solution}

\begin{problem}
已知 \(a \in \R\)，函数 \(f(x) = x^2 |x - a|\).
\begin{enumerate}
\item 当 \(a = 2\) 时，求 \(f(x) = x\) 使成立的 \(x\) 的集合；
\item 求函数 \(y = f(x)\) 在区间 \([1,2]\) 上的最小值.
\end{enumerate}
\end{problem}

\begin{solution}
\begin{enumerate}
\item 当 \(x<2\) 时，\(|x-2|=2-x\)，方程化为

\(x^2(2-x)=x\iff x\left[x(2-x)-1\right]=0\iff -x(x-1)^2=0\).

所以得到 \(x=0\) 或 \(x=1\)，这两个解都满足 \(x<2\).

当 \(x\geqslant2\) 时，\(|x-2|=x-2\)，方程化为 \(x^2(x-2)=x\). 因为此时 \(x>0\)，可除以 \(x\)，得到 \(x^2-2x-1=0\)，所以 \(x=1\pm\sqrt{2}\). 其中只有 \(1+\sqrt{2}\geqslant2\)，故所求解集为 \(\{0,1,1+\sqrt{2}\}\).
\item 记 \([1,2]\) 上的最小值为 \(m\). 由于 \(x\in[1,2]\)，按参数 \(a\) 的位置分类.

当 \(a\leqslant1\) 时，\(x-a\geqslant0\)，所以 \(f(x)=x^3-ax^2\)，且

\(f'(x)=x(3x-2a)\geqslant1\times(3-2a)>0\).

因此 \(f\) 在 \([1,2]\) 上递增，最小值为 \(f(1)=1-a\).

当 \(1<a\leqslant2\) 时，点 \(x=a\) 属于 \([1,2]\)，且 \(f(a)=a^2|a-a|=0\). 又 \(f(x)\geqslant0\)，故最小值为 \(0\).

当 \(a>2\) 时，\(x-a<0\)，所以 \(f(x)=ax^2-x^3\)，且

\(f'(x)=x(2a-3x)\).

若 \(a\geqslant3\)，则对 \(x\in[1,2]\) 有 \(2a-3x\geqslant2a-6\geqslant0\)，所以 \(f\) 递增，最小值为 \(f(1)=a-1\).

若 \(2<a<3\)，则 \(\dfrac{2a}{3}\in(1,2)\)，函数在 \(\left[1,\dfrac{2a}{3}\right]\) 上递增，在 \(\left[\dfrac{2a}{3},2\right]\) 上递减. 因此最小值只能在两个端点处取得. 计算

\(f(1)=a-1\)，\(f(2)=4(a-2)\).

比较两者，\(4(a-2)\leqslant a-1\iff a\leqslant\dfrac73\). 因此在 \(2<a<3\) 时，\(a\leqslant\dfrac73\) 取 \(f(2)\)，而 \(a>\dfrac73\) 取 \(f(1)\). 与前面的情形合并，所求最小值为
\examdisplaycases{m =}{
1-a, & a\leqslant 1,\\
0, & 1<a\leqslant 2,\\
4(a-2), & 2<a\leqslant \dfrac{7}{3},\\
a-1, & a>\dfrac{7}{3}.
}
\end{enumerate}
\end{solution}

\begin{problem}
设数列 \(\{a_{n}\}\) 的前 \(n\) 项和为 \(S_{n}\)，已知 \(a_1 = 1\)，\(a_2 = 6\)，\(a_3 = 11\)，且 \((5n - 8)S_{n+1} - (5n + 2)S_n = An + B\)，\(n = 1\)，\(2\)，\(3\)，\(\cdots\)，其中 \(A\)，\(B\) 为常数.
\begin{enumerate}
\item 求 \(A\) 与 \(B\) 的值；
\item 证明数列 \(\{a_{n}\}\) 为等差数列；
\item 证明不等式 \(\sqrt{5a_{mn}} - \sqrt{a_m a_n} > 1\) 对任何正整数 \(m\)，\(n\) 都成立.
\end{enumerate}
\end{problem}

\begin{solution}
\begin{enumerate}
\item 由已知首三项，

\(S_1=a_1=1\)，\(S_2=a_1+a_2=7\)，\(S_3=a_1+a_2+a_3=18\).

在题设关系中取 \(n=1\)，得到

\((5-8)S_2-(5+2)S_1=A+B\iff-3\times7-7\times1=A+B\)，

即 \(A+B=-28\). 取 \(n=2\)，得到

\((10-8)S_3-(10+2)S_2=2A+B\iff2\times18-12\times7=2A+B\)，

即 \(2A+B=-48\). 两式相减得 \(A=-20\)，再代回得 \(B=-8\).
\item 由第一问，题设关系化为

\((5n-8)S_{n+1}-(5n+2)S_n=-20n-8\).

把其中的 \(n\) 换成 \(n+1\)，得到

\((5n-3)S_{n+2}-(5n+7)S_{n+1}=-20n-28\).

用这个等式减去原等式，消去右端的含 \(n\) 一次项，得

\((5n-3)S_{n+2}-(10n-1)S_{n+1}+(5n+2)S_n=-20\).

再把这个等式中的 \(n\) 换成 \(n+1\)，并减去上式，得到

\((5n+2)S_{n+3}-(15n+6)S_{n+2}+(15n+6)S_{n+1}-(5n+2)S_n=0\).

注意到 \(15n+6=3(5n+2)\)，并用

\(S_{n+1}=S_n+a_{n+1}\)，\(S_{n+2}=S_n+a_{n+1}+a_{n+2}\)，\(S_{n+3}=S_n+a_{n+1}+a_{n+2}+a_{n+3}\)，

上式中关于 \(S_n\) 的项相互抵消，剩下

\((5n+2)(a_{n+3}-2a_{n+2}+a_{n+1})=0\).

由于 \(n\geqslant1\) 时 \(5n+2\neq0\)，所以

\(a_{n+3}-a_{n+2}=a_{n+2}-a_{n+1}\).

这说明从第二项开始相邻两项的差相等. 又 \(a_2-a_1=5\)，\(a_3-a_2=5\)，故数列 \(\{a_n\}\) 为公差 \(5\) 的等差数列，通项为

\(a_n=a_1+(n-1)\times5=5n-4\).
\item 由上一问，

\(a_{mn}=5mn-4\)，\(a_m=5m-4\)，\(a_n=5n-4\).

令 \(X=5a_{mn}=25mn-20\)，\(Y=a_ma_n=(5m-4)(5n-4)=25mn-20m-20n+16\). 因为 \(m\)，\(n\) 是正整数，所以 \(a_m\)，\(a_n>0\)，从而 \(Y>0\). 直接计算得

\(X-Y-1=20m+20n-37\).

由基本不等式，

\(2\sqrt{Y}=2\sqrt{a_ma_n}\leqslant a_m+a_n=5m+5n-8\).

又因为 \(m\)，\(n\geqslant1\)，

\((20m+20n-37)-(5m+5n-8)=15m+15n-29\geqslant1>0\)，

所以 \(2\sqrt{Y}<20m+20n-37=X-Y-1\). 因而

\(X>Y+1+2\sqrt{Y}=(\sqrt{Y}+1)^2\).

由于 \(X>0\) 且 \(Y>0\)，两边开平方得 \(\sqrt{X}>\sqrt{Y}+1\)，即

\(\sqrt{5a_{mn}}-\sqrt{a_ma_n}>1\).

这对任意正整数 \(m\)，\(n\) 均成立.
\end{enumerate}
\end{solution}
